\documentclass[12pt,a4paper]{article}

% ===========================
% PACKAGES - Available subset
% ===========================
\usepackage[utf8]{inputenc}
\usepackage[T1]{fontenc}
\usepackage[brazilian]{babel}

\usepackage[
  left=2.5cm,
  right=2.5cm,
  top=3cm,
  bottom=3cm
]{geometry}

\usepackage{amsmath,amssymb,amsthm,mathtools}
\usepackage{bm}
\usepackage{stmaryrd}
\usepackage{proof}

\usepackage{tikz}
\usepackage{tikz-cd}
\usetikzlibrary{positioning,arrows.meta,shapes}

\usepackage{listings}
\usepackage{booktabs,array,multirow}
\usepackage{fancyhdr}
\usepackage{lastpage}
\usepackage{hyperref}
\usepackage{xcolor}
\usepackage{mdframed}
\usepackage{tcolorbox}
\tcbuselibrary{theorems,skins,breakable}
\usepackage{enumitem}

% ===========================
% CORES (from itt-lean.cls)
% ===========================
\definecolor{leanblue}{RGB}{0, 84, 166}
\definecolor{leangreen}{RGB}{0, 128, 0}
\definecolor{leanred}{RGB}{220, 20, 60}
\definecolor{leanorange}{RGB}{255, 140, 0}
\definecolor{leanpurple}{RGB}{128, 0, 128}
\definecolor{proofgray}{RGB}{245, 245, 245}

% ===========================
% HYPERREF
% ===========================
\hypersetup{
  colorlinks=true,
  linkcolor=leanblue,
  citecolor=leangreen,
  urlcolor=leanpurple,
}

% ===========================
% INFORMAÇÕES DO DOCUMENTO
% ===========================
\newcommand{\nomealuno}{Igor Strozzi}
\newcommand{\matricula}{2025123456}
\newcommand{\disciplina}{Teoria dos Tipos Intuicionista e Lean}
\newcommand{\professor}{Francesco Noseda}
\newcommand{\semestre}{2025-2}
\newcommand{\lista}{S}
\newcommand{\universidade}{UFRJ}
\newcommand{\departamento}{Instituto de Matemática}

% ===========================
% CABEÇALHO E RODAPÉ
% ===========================
\pagestyle{fancy}
\fancyhf{}
\lhead{\nomealuno}
\chead{\disciplina}
\rhead{Lista \lista}
\lfoot{\universidade}
\cfoot{}
\rfoot{Página \thepage\ de \pageref*{LastPage}}
\renewcommand{\headrulewidth}{0.4pt}
\renewcommand{\footrulewidth}{0.4pt}

% ===========================
% PÁGINA DE TÍTULO
% ===========================
\newcommand{\maketitlecustom}{
  \thispagestyle{empty}
  \begin{center}
    {\Large \textbf{\universidade}} \\[0.3cm]
    {\large \departamento} \\[0.5cm]
    \hrule
    \vspace{0.4cm}
    {\LARGE \textbf{\disciplina}} \\[0.3cm]
    {\Large Lista de Exercícios \lista} \\[0.5cm]
    \hrule
    \vspace{0.5cm}
    \begin{tabular}{ll}
      \textbf{Aluno:} & \nomealuno \\
      \textbf{Matrícula:} & \matricula \\
      \textbf{Professor:} & \professor \\
      \textbf{Semestre:} & \semestre \\
      \textbf{Data:} & \today
    \end{tabular}
  \end{center}
  \vspace{1cm}
}

% ===========================
% AMBIENTES MATEMÁTICOS
% ===========================
\theoremstyle{definition}
\newtheorem{definicao}{Definição}[section]
\newtheorem{exemplo}[definicao]{Exemplo}

\theoremstyle{plain}
\newtheorem{teorema}[definicao]{Teorema}
\newtheorem{proposicao}[definicao]{Proposição}
\newtheorem{lema}[definicao]{Lema}
\newtheorem{corolario}[definicao]{Corolário}

\theoremstyle{remark}
\newtheorem{observacao}[definicao]{Observação}

% Caixa colorida para teoremas
\tcolorboxenvironment{teorema}{
  boxrule=0pt,
  boxsep=0pt,
  colback={leanblue!10},
  left=2mm,
  right=2mm,
  top=2mm,
  bottom=2mm,
  before skip=10pt,
  after skip=10pt,
  borderline west={2pt}{0pt}{leanblue}
}

% ===========================
% COMANDOS PARA LÓGICA
% ===========================
\newcommand{\e}{\wedge}
\newcommand{\ou}{\vee}
\newcommand{\implic}{\to}
\newcommand{\nao}{\neg}
\newcommand{\falsum}{\bot}
\newcommand{\verum}{\top}
\newcommand{\deriva}{\vdash}
\newcommand{\satisfaz}{\models}
\newcommand{\sem}[1]{\llbracket #1 \rrbracket}
\newcommand{\defeq}{\coloneqq}

% ===========================
% LEAN CODE
% ===========================
\lstdefinelanguage{Lean}{
  keywords={theorem, lemma, def, example, by, intro, exact, apply, rfl, simp, sorry, constructor, cases, induction, have, show, assume, fun, match, with, end, if, then, else, let, in, where, structure, class, instance, extends, namespace, variable, open, notation, inductive, deriving, calc, universe, Type, Prop},
  sensitive=true,
  comment=[l]{--},
  morecomment=[s]{/-}{-/},
  string=[b]",
}

\lstnewenvironment{leancode}{%
  \lstset{
    language=Lean,
    basicstyle=\ttfamily\small,
    keywordstyle=\color{leanblue}\bfseries,
    commentstyle=\color{leangreen}\itshape,
    stringstyle=\color{leanred},
    breaklines=true,
    frame=single,
    rulecolor=\color{leanblue},
    backgroundcolor=\color{proofgray},
    xleftmargin=1em,
    xrightmargin=1em,
  }%
}{}

\newcommand{\lean}[1]{\texttt{\color{leanblue}#1}}

% ===========================
% RETICULADOS E HEYTING
% ===========================
\tikzset{
  lattice node/.style={circle, draw, fill=white, inner sep=2pt, minimum size=6mm},
  lattice edge/.style={draw, thick},
  heyting node/.style={circle, draw, fill=white, inner sep=2pt, minimum size=6mm}
}

\newcommand{\heytingtres}{%
  \begin{tikzpicture}[scale=0.8]
    \node[heyting node] (0) at (0,0) {$\bot$};
    \node[heyting node] (a) at (0,1.5) {$a$};
    \node[heyting node] (1) at (0,3) {$\top$};
    \draw[lattice edge] (0) -- (a) -- (1);
  \end{tikzpicture}%
}

% ===========================
% DOCUMENTO
% ===========================
\begin{document}
\maketitlecustom

\begin{center}
  \Large\textbf{Ontologias Mereotopológicas com Álgebras Composicionais:\\
  Uma Formalização em Lean 4}
\end{center}

\vspace{1em}

\tableofcontents
\newpage

\section{Introdução}

Este documento desenvolve os fundamentos matemáticos para formalizar \emph{ontologias mereotopológicas} equipadas com \emph{álgebras composicionais} que funcionam como sistemas dedutivos. A percepção central é que as álgebras de Heyting---a semântica algébrica da lógica intuicionista---fornecem o arcabouço natural para tais estruturas composicionais.

Dada uma ontologia formal $\mathcal{O}$ com estrutura mereotopológica (relações parte-todo e conexão topológica) e uma álgebra $\mathcal{A}$ agindo sobre $\mathcal{O}$, buscamos formalizar:
\begin{enumerate}
  \item A estrutura mereotopológica de $\mathcal{O}$
  \item A estrutura algébrica de $\mathcal{A}$ como álgebra de Heyting
  \item A ação $\sem{\cdot} : \mathcal{O} \to A$ preservando estrutura relevante
  \item O caráter dedutivo emergente da implicação de Heyting
\end{enumerate}

\section{Ontologia Mereotopológica}

\begin{definicao}[Ontologia Mereotopológica]
Uma \emph{ontologia mereotopológica} é uma estrutura $\mathcal{O} = (O, \sqsubseteq, \mathsf{C})$ onde:
\begin{itemize}
  \item $O$ é um conjunto não-vazio de \emph{conceitos} ou \emph{entidades}
  \item $\sqsubseteq$ é uma relação de \emph{partição} (ordem parcial)
  \item $\mathsf{C}$ é uma relação de \emph{conexão} (reflexiva, simétrica)
\end{itemize}
satisfazendo o \emph{axioma de coerência mereotopológica}:
$$
x \sqsubseteq y \implies \mathsf{C}(x, y)
$$
\end{definicao}

A relação de partição $\sqsubseteq$ satisfaz:
\begin{itemize}
  \item \textbf{Reflexividade:} $\forall x.\, x \sqsubseteq x$
  \item \textbf{Transitividade:} $x \sqsubseteq y \land y \sqsubseteq z \implies x \sqsubseteq z$
  \item \textbf{Antissimetria:} $x \sqsubseteq y \land y \sqsubseteq x \implies x = y$
\end{itemize}

\begin{definicao}[Relações Mereotopológicas Derivadas]
Das relações primitivas, definimos:
\begin{align*}
  \mathsf{PP}(x, y) &\defeq x \sqsubseteq y \land x \neq y && \text{(parte própria)} \\
  \mathsf{O}(x, y) &\defeq \exists z.\, z \sqsubseteq x \land z \sqsubseteq y && \text{(sobreposição)} \\
  \mathsf{D}(x, y) &\defeq \neg \mathsf{O}(x, y) && \text{(disjunção)} \\
  \mathsf{EC}(x, y) &\defeq \mathsf{C}(x, y) \land \mathsf{D}(x, y) && \text{(conexão externa)}
\end{align*}
\end{definicao}

\begin{observacao}
A relação de conexão $\mathsf{C}$ captura adjacência topológica sem exigir partes compartilhadas. Dois conceitos podem estar externamente conectados (tocando-se) sem sobreposição---uma distinção crucial para raciocínio espacial e conceitual.
\end{observacao}

\section{Álgebras de Heyting como Álgebras Composicionais}

Recordemos do Capítulo 2 (Definição 2.4.2) que uma álgebra de Heyting fornece semântica algébrica para a lógica intuicionista.

\begin{definicao}[Álgebra de Heyting]
Uma \emph{álgebra de Heyting} é um reticulado distributivo limitado $\mathcal{H} = (H, \sqcup, \sqcap, \Rightarrow, 0, 1)$ tal que para todos $a, b \in H$, o \emph{pseudo-complemento relativo} $a \Rightarrow b$ existe, caracterizado pela \textbf{propriedade de adjunção}:
$$
c \sqcap a \leq b \iff c \leq a \Rightarrow b
$$
\end{definicao}

\begin{teorema}[Completude]
Para a lógica proposicional intuicionista (Teorema 2.4.7):
$$
\Gamma \vdash \varphi \iff \Gamma \models \varphi
$$
onde $\Gamma \models \varphi$ significa: para toda álgebra de Heyting $\mathcal{H}$ e valoração $v$, se $\mathcal{H}, v \models \Gamma$ então $\mathcal{H}, v \models \varphi$.
\end{teorema}

\subsection{Estrutura da Álgebra Composicional}

Interpretamos as operações da álgebra de Heyting como operadores de composição conceitual:

\begin{center}
\begin{tabular}{cll}
  \textbf{Operação} & \textbf{Notação} & \textbf{Interpretação} \\
  \midrule
  Meet & $a \sqcap b$ & Conjunção/composição conceitual \\
  Join & $a \sqcup b$ & Disjunção/escolha conceitual \\
  Implicação & $a \Rightarrow b$ & Acarretamento conceitual \\
  Negação & $\sim a = a \Rightarrow 0$ & Complemento conceitual \\
  Topo & $1$ & Conceito universal \\
  Fundo & $0$ & Conceito vazio/impossível \\
\end{tabular}
\end{center}

\begin{definicao}[Derivabilidade]
Em uma álgebra composicional $\mathcal{A}$, definimos a \emph{relação de derivabilidade}:
$$
a \vdash_{\mathcal{A}} b \defeq a \sqcap b = a
$$
Isto é equivalente a $a \leq b$ na ordem subjacente.
\end{definicao}

\section{A Álgebra Ontológica}

\begin{definicao}[Álgebra Ontológica]
Uma \emph{álgebra ontológica} é uma tripla $(\mathcal{O}, \mathcal{A}, \sem{\cdot})$ onde:
\begin{itemize}
  \item $\mathcal{O} = (O, \sqsubseteq, \mathsf{C})$ é uma ontologia mereotopológica
  \item $\mathcal{A}$ é uma álgebra de Heyting (álgebra composicional)
  \item $\sem{\cdot} : O \to A$ é uma \emph{função de interpretação}
\end{itemize}
satisfazendo:
\begin{enumerate}
  \item \textbf{Monotonicidade:} $c_1 \sqsubseteq c_2 \implies \sem{c_1} \leq \sem{c_2}$
  \item \textbf{Não-trivialidade de conexão:} $\mathsf{C}(c_1, c_2) \implies \sem{c_1} \sqcap \sem{c_2} \neq 0$
\end{enumerate}
\end{definicao}

\begin{definicao}[Operações Fundamentais]
Dados conceitos $C_1, C_2 \in O$, definimos:
\begin{align*}
  C_1 \triangleright C_2 &\defeq \sem{C_1} \sqcap \sem{C_2} && \text{(combinação)} \\
  C_1 \multimap C_2 &\defeq \sem{C_1} \Rightarrow \sem{C_2} && \text{(implicação conceitual)}
\end{align*}
\end{definicao}

\subsection{Teoremas Principais}

\begin{teorema}[Partição Implica Derivabilidade]
Se $c_1 \sqsubseteq c_2$ em $\mathcal{O}$, então:
$$
\sem{c_1} \sqcap \sem{c_2} = \sem{c_1}
$$
Isto é, $\sem{c_1} \vdash_{\mathcal{A}} \sem{c_2}$.
\end{teorema}

\begin{proof}
Por monotonicidade, $\sem{c_1} \leq \sem{c_2}$. Então:
$$
\sem{c_1} \sqcap \sem{c_2} \leq \sem{c_1} \leq \sem{c_2}
$$
logo $\sem{c_1} \sqcap \sem{c_2} \leq \sem{c_1}$. Reciprocamente, $\sem{c_1} \leq \sem{c_1} \sqcap \sem{c_2}$ pois $\sem{c_1} \leq \sem{c_1}$ e $\sem{c_1} \leq \sem{c_2}$.
\end{proof}

\begin{teorema}[Não-Trivialidade de Sobreposição]
Se $\mathsf{O}(c_1, c_2)$ em $\mathcal{O}$, então $c_1 \triangleright c_2 \neq 0$.
\end{teorema}

\begin{teorema}[Modus Ponens para Conceitos]
Para quaisquer $c_1, c_2 \in O$:
$$
\sem{c_1} \sqcap (c_1 \multimap c_2) \leq \sem{c_2}
$$
\end{teorema}

\section{Conexão com Semântica Topológica}

Conjuntos abertos de um espaço topológico formam uma álgebra de Heyting (Proposição 2.4.8).

\begin{proposicao}[Álgebra de Heyting Topológica]
Seja $\mathcal{T}$ um espaço topológico. Então $\mathcal{H} = (\mathcal{O}(\mathcal{T}), \cup, \cap, \Rightarrow, \emptyset, \mathcal{T})$ é uma álgebra de Heyting onde:
\begin{align*}
  A \Rightarrow B &= \mathrm{Int}(-A \cup B) \\
  \sim A &= \mathrm{Int}(-A)
\end{align*}
\end{proposicao}

\begin{exemplo}[Falha da Lei de Peirce]
Considere $\mathcal{H} = \mathcal{O}(\mathbb{R})$ com $v(p) = \mathbb{R} \setminus \{0\}$ e $v(q) = \emptyset$. Então:
$$
\sem{(p \to q) \to p}_v = \mathrm{Int}(\emptyset \cup (\mathbb{R} \setminus \{0\})) = \mathbb{R} \setminus \{0\}
$$
Logo $\sem{((p \to q) \to p) \to p}_v = \mathbb{R} \setminus \{0\} \neq \mathbb{R}$.

\textbf{Ontologicamente:} A lei de Peirce falha porque não podemos construtivamente obter $p$ meramente de saber que ``se $p$ implicasse $q$, então $p$ seria válido.''
\end{exemplo}

\section{O Exemplo de Três Elementos}

Considere a cadeia de três elementos $\mathbf{3} = \{0 < a < 1\}$:

\begin{center}
\heytingtres
\end{center}

As operações são:
\begin{center}
\begin{tabular}{c|ccc}
  $\sqcap$ & 0 & $a$ & 1 \\
  \hline
  0 & 0 & 0 & 0 \\
  $a$ & 0 & $a$ & $a$ \\
  1 & 0 & $a$ & 1
\end{tabular}
\qquad
\begin{tabular}{c|ccc}
  $\Rightarrow$ & 0 & $a$ & 1 \\
  \hline
  0 & 1 & 1 & 1 \\
  $a$ & 0 & 1 & 1 \\
  1 & 0 & $a$ & 1
\end{tabular}
\end{center}

\begin{observacao}
Note que $\sim a = a \Rightarrow 0 = 0$, mas $\sim\sim a = 0 \Rightarrow 0 = 1 \neq a$. Isto demonstra a falha da eliminação da dupla negação na lógica intuicionista (cf. Exemplo 2.1.4(iii)).
\end{observacao}

\section{Formalização em Lean 4}

A formalização completa é autocontida, construindo todas as estruturas a partir de primeiros princípios.

\subsection{Fundamentos Algébricos}

\begin{leancode}
class PartialOrder' (α : Type u) where
  le : α → α → Prop
  le_refl : ∀ a, le a a
  le_trans : ∀ {a b c}, le a b → le b c → le a c
  le_antisymm : ∀ {a b}, le a b → le b a → a = b

class HeytingAlgebra (α : Type u) extends BoundedLattice α where
  himp : α → α → α
  le_himp_iff : ∀ {a b c}, le (inf a b) c ↔ le a (himp b c)
\end{leancode}

\subsection{Ontologia Mereotopológica}

\begin{leancode}
class MereotopologicalOntology (O : Type u) where
  partOf : O → O → Prop
  C : O → O → Prop
  partOf_refl : ∀ x, partOf x x
  partOf_trans : ∀ {x y z}, partOf x y → partOf y z → partOf x z
  partOf_antisymm : ∀ {x y}, partOf x y → partOf y x → x = y
  C_refl : ∀ x, C x x
  C_symm : ∀ {x y}, C x y → C y x
  part_connected : ∀ {x y}, partOf x y → C x y
\end{leancode}

\subsection{A Álgebra Ontológica}

\begin{leancode}
class OntologicalAlgebra (O : Type u) (A : Type v)
    [MereotopologicalOntology O] [CompositionalAlgebra A] where
  interp : O → A
  interp_mono : ∀ {c₁ c₂}, partOf c₁ c₂ → interp c₁ ≤ interp c₂
  connected_nontrivial : ∀ {c₁ c₂}, C c₁ c₂ → 
    inf (interp c₁) (interp c₂) ≠ bot
\end{leancode}

\subsection{Teorema Principal em Lean}

\begin{leancode}
theorem partOf_derives {c₁ c₂ : O}
    (h : MereotopologicalOntology.partOf c₁ c₂) :
    ⦃c₁⦄ ⊗ ⦃c₂⦄ = ⦃c₁⦄ := by
  apply le_antisymm
  · exact inf_le_left ⦃c₁⦄ ⦃c₂⦄
  · exact le_inf (le_refl ⦃c₁⦄) (interp_mono h)
\end{leancode}

\section{Observações Filosóficas}

A conexão entre ontologia mereotopológica e álgebras de Heyting ilumina correspondências profundas:

\begin{enumerate}
  \item \textbf{Ontologia Construtiva:} Assim como a lógica intuicionista requer evidência construtiva para asserções, a composição ontológica requer testemunhas explícitas.
  
  \item \textbf{Falha da Dualidade Clássica:} A assimetria entre $\varphi \to \neg\neg\varphi$ (válida) e $\neg\neg\varphi \to \varphi$ (não válida) espelha a assimetria em mereotopologia.
  
  \item \textbf{Intuição Topológica:} Interpretar conceitos como conjuntos abertos fornece intuição geométrica.
  
  \item \textbf{Semântica Composicional:} A adjunção $a \sqcap b \leq c \Leftrightarrow a \leq b \Rightarrow c$ captura a essência do raciocínio composicional.
\end{enumerate}

\section{Conclusão}

Estabelecemos um arcabouço rigoroso conectando mereotopologia, álgebras de Heyting, raciocínio composicional e formalização em Lean 4. A estrutura de álgebra de Heyting captura naturalmente o caráter intuicionista do raciocínio ontológico.

\section*{Referências}

\begin{enumerate}
  \item Sørensen, M.H. \& Urzyczyn, P. \textit{Lectures on the Curry-Howard Isomorphism}. Capítulo 2: Lógica Intuicionista.
  \item Casati, R. \& Varzi, A.C. \textit{Parts and Places: The Structures of Spatial Representation}. MIT Press, 1999.
  \item Heyting, A. \textit{Intuitionism: An Introduction}. North-Holland, 1956.
\end{enumerate}

\end{document}
